\chapter{从现实问题到机器学习问题：建模的艺术}

AI 最大的瓶颈往往不在算力，而在于使用者能否把问题说清楚。本章讨论如何将模糊的现实问题转化为可计算的机器学习问题，并将建模决策通过 Prompt 传递给 AI。

\section{为什么问题建模是人类不可替代的核心能力}

现实任务通常没有天然明确的目标、标签和评价标准。AI 可以根据给定定义执行建模，但无法替使用者决定什么问题值得解决，以及什么结果才算正确。

\section{现实问题到机器学习问题的翻译框架}

\begin{enumerate}
  \item \textbf{对象是谁}：道路、交叉口、OD 对、乘客、司机、车辆、公交站点或城市区域。
  \item \textbf{目标是什么}：预测、解释、分类、发现结构、优化决策或评估政策效果。
  \item \textbf{标签是什么}：例如“拥堵”使用速度阈值还是延误定义，“事故严重程度”使用警方记录、医疗结果还是财产损失。
  \item \textbf{时间和空间尺度是什么}：5 分钟、15 分钟、小时或日；路段、区域或城市。
  \item \textbf{错误成本是什么}：事故漏报、拥堵误判、调度失败和补贴超预算的代价完全不同。
\end{enumerate}

\section{机器学习问题的主要学习范式}

\subsection{监督学习}

\begin{itemize}
  \item 回归：预测连续值，例如交通流量预测和出行时间估计。
  \item 分类：判断所属类别，例如事故严重程度分级和驾驶行为识别。
\end{itemize}

\subsection{无监督学习}

\begin{itemize}
  \item 聚类：发现隐藏结构，例如出行模式挖掘和路段相似性分组。
  \item 降维：理解高维数据结构。
  \item 异常检测：发现不寻常状态。
\end{itemize}

\subsection{强化学习}

强化学习面向序列决策问题，在状态、动作和奖励之间学习策略，例如信号控制和动态定价。

\subsection{交通数据的特殊结构}

\begin{itemize}
  \item 时序结构：交通流、客流和需求随时间变化。
  \item 空间结构：路段、站点和区域之间存在关联。
  \item 图结构：道路网络、公交网络和 OD 网络天然可以表示为图。
\end{itemize}

\section{领域知识的注入：如何把专业知识变成Prompt的一部分}

需要明确哪些信息是 AI 不知道、但领域专家知道的，例如数据采集机制、字段定义边界和业务规则。可以使用如下表达：

\begin{quote}
在我的数据中，字段 X 的含义是……，与标准定义不同之处在于……，请在分析时注意这一点。
\end{quote}

\section{实战练习：现实问题的建模全过程}

以“缓解校园周边拥堵”为例，它可以被转化为多种任务：

\begin{enumerate}
  \item 预测任务：预测校门口 15 分钟后的车速。
  \item 分类任务：判断是否进入拥堵状态。
  \item 优化任务：调整信号灯或停车管理策略。
  \item 因果任务：评估限行措施是否真的降低拥堵。
  \item 生成式任务：让 AI 生成交通组织方案，但必须人工审查。
\end{enumerate}
